\documentclass[12pt]{article}

\usepackage[utf8]{inputenc}
\usepackage{latexsym,amsfonts,amssymb,amsthm,amsmath}
\usepackage{tikz}
\usepackage{stanli}
\usetikzlibrary{shapes,calc,through,3d,patterns}
\usetikzlibrary{arrows.meta,decorations.pathreplacing,angles,quotes}
\usepackage{multicol}
\usepackage{geometry}
\usepackage{mathtools}
\usepackage{siunitx}
\usepackage{booktabs}

\setlength{\parindent}{0in}
\setlength{\oddsidemargin}{-0.25in}
\setlength{\textwidth}{7in}
\setlength{\textheight}{9.3in}
\setlength{\topmargin}{-1in}
\setlength{\headheight}{18pt}

\newcommand\blfootnote[1]{%
    \begingroup
    \renewcommand\thefootnote{}\footnote{#1}%
    \addtocounter{footnote}{-1}%
    \endgroup
}

\definecolor{darkred}{rgb}{0.6,0,0}

\newcommand{\bolddelta}{\boldsymbol\delta}
\newcommand{\boldxi}{\boldsymbol\xi}

\newcommand{\T}{^{T}}
\newcommand{\bK}{\overline{\mathbf K}}
\newcommand{\bM}{\overline{\mathbf M}}
\newcommand{\vct}[1]{\mathbf{#1}}
\newcommand{\colv}[1]{\begin{bmatrix}#1\end{bmatrix}}


\begin{document}

\noindent ME473 - Dynamic finite element analysis of structures\hfill EPFL\\
Stefano Burzio \hfill 2025\\
\vspace{0.3cm}

\hspace*{\fill}\textbf{\Large Exam questions}\hspace*{\fill}

\hrulefill

\subsection*{Question 1}
Let
        \[
          \vct K=\colv{3&1\\1&7},\qquad \vct M=\colv{1&1\\1&4}.
        \]
        First apply the standard angle in Jacobi method chosen from $\vct K$ and check the mass
        off-diagonal entry. Then compute $x,\alpha,\gamma$ using the generalieed Jacobi method.
     Reads  the two generalised eigenvalues from the diagonal entries of the rotated matrices. 

For the given stiffness block, $\tan(2\theta)=-1/2$.
The small-angle branch has $\cos(2\theta)=2/\sqrt5$ and
$\sin(2\theta)=-1/\sqrt5$. It clears the stiffness pivot, but the rotated
mass still has the off-diagonal entry
\[
   m'_{ij}=m_{ij}\cos(2\theta)-\tfrac12(m_{ii}-m_{jj})\sin(2\theta)
   =\frac{2}{\sqrt5}-\frac{3}{2\sqrt5}=\frac{1}{2\sqrt5}.
\]
For the generalized transformation the intermediate quantities are
\[
   \bar k_{ii}=3(1)-1(1)=2,\qquad
   \bar k_{jj}=7(1)-4(1)=3,\qquad
   \bar k=\tfrac12[3(4)-7(1)]=\tfrac52 .
\]
The square root is $7/2$, so $x=6$ and
\[
   \boxed{\alpha=\tfrac12,\qquad\gamma=-\tfrac13},\qquad
   \vct J=\colv{1&1/2\\-1/3&1}.
\]
Here $\det\vct J=7/6$. Direct multiplication gives
\[
   \vct J\T\vct K_p\vct J=\colv{28/9&0\\0&35/4},\qquad
   \vct J\T\vct M_p\vct J=\colv{7/9&0\\0&21/4}.
\]
The diagonal ratios are $4$ and $5/3$, respectively. The corresponding
columns of $\vct J$ are eigenvectors of the original pencil. As an
independent check, $\det(\vct K_p-\lambda\vct M_p)
   =3\lambda^2-17\lambda+20=(\lambda-4)(3\lambda-5)$.
   
\end{document}
